% r2_section.tex — SKELETON (crew flavor-obstruction-paper, R2, 2026-10-23)
% R2: Chazz's open mandate ("what fixes theta23, by any route"), the manager's brief
% galt-crews/flavor-obstruction/files/brief-R2.md. A fragment for \input after
% constructive_section.tex. It uses the macros of theta23_obstruction_paper.tex and touches no
% proof there. Every number is [pending] until the lab's record for that candidate lands and is
% confirmed through the manager. Inputs: analysis/data/nufit6_ranges.csv (728a1d0) and
% analysis/data/r2_inputs.csv (pending, literature item #12).

%=======================================================================
\section{Beyond the relational class: other sources of orientation}\label{sec:r2}
%=======================================================================
The no-go and the battery of Sec.~\ref{sec:constructive} show that nothing built from $(\He,\Hnu)$,
and no added flavor dimension, orients the $\nu_2$--$\nu_3$ frame. This section tests sources that
are \emph{other data} or \emph{input-free}, together with the hypothesis that nothing needs
explaining. The criteria are (a)--(e) of Sec.~\ref{sec:conditions}, with (b) widened: oriented by
data other than $\theta_{23}$, \emph{or by no continuous input at all}. Every candidate is scored
on two questions: does it give near-maximal mixing, and which octant does it pick?

\subsection{What is established: near-maximal, not the octant}\label{sec:octant}
% PENDING: theorist item #13 (the paragraph, with the sigma distances computed from
% analysis/data/nufit6_ranges.csv). Brief values to verify: 45 deg is 3.9 sigma from IC19 w/o
% SK-atm NO (48.5 +0.7/-0.9) and 1.7 sigma from IC24 with SK-atm NO (43.3 +1.0/-0.8).
[pending]

\subsection{D7: residual-symmetry sum rules}\label{sec:D7}
\emph{Candidate.} TM1, TM2 and $\mu$--$\tau$ reflection, which give closed-form relations between
$\theta_{23}$, $\theta_{13}$, $\delta$ (and $\theta_{12}$), plus the general first-order rules
$\theta_{23}-45^\circ\sim\theta_{13}\cos\delta$.\\
\emph{Gate.} Reproduce each rule against its source (exact forms: item \#12).\\
\emph{Question.} With $\theta_{13}$ and $\delta$ at NuFIT~6.0 (central values and $3\sigma$
$\delta$ ranges, all four variants), what $\theta_{23}$ and octant does each rule give, and is the
result inside the $3\sigma$ window? \emph{Null.} The window is about $10^\circ$ wide, so the score
is the joint prediction ($\theta_{12}$ too, where the rule fixes it) and the sign of the
$\delta$--octant correlation. (b) passes by construction; the test is (e).\\
\emph{Result.} [pending]\quad\emph{Verdict.} [pending]

\subsection{D8: quark--lepton complementarity}\label{sec:D8}
\emph{Candidate.} $U_{\rm PMNS}=U_{\rm bimax}V_{\rm CKM}^\dagger$ (and the other order):
$\theta_{12}+\theta_C\simeq45^\circ$, $\theta_{23}+\theta_{23}^{\rm CKM}\simeq45^\circ$.\\
\emph{Gate.} Reproduce Minakata--Smirnov and Raidal (2004) with PDG CKM inputs (item \#12).\\
\emph{Question.} The exact $\theta_{23}$ (and $\theta_{12}$, $\theta_{13}$, $\delta$) for each
composition order and phase convention, against NuFIT; the octant each order picks; whether any
order gets $\theta_{13}$ right. \emph{Null.} Random correction matrices with the CKM's angle
sizes.\\
\emph{Result.} [pending]\quad\emph{Verdict.} [pending]

\subsection{D9: Haar typicality (anarchy)}\label{sec:D9}
\emph{Candidate.} No mechanism: the mixing matrix is a typical Haar draw on $U(3)$.\\
\emph{Gate.} Reproduce de~Gouv\^ea--Murayama: the flat Haar-induced variables and their reported
p-value (item \#12).\\
\emph{Question.} The p-value of the NuFIT~6.0 point under Haar, and the Bayes factor of Haar
against a $\mu$--$\tau$-reflection-centred prior (width from the D2 breaking size). If Haar is not
disfavoured, the verdict is that no orientation mechanism is required by the data.\\
\emph{Result.} [pending]\quad\emph{Verdict.} [pending]

\subsection{D10: modular symmetry at fixed points}\label{sec:D10}
\emph{Candidate.} An $A_4$ modular model at the fixed points $\tau=i,\ \omega,\ i\infty$
(residual $\mathbb Z_2$, $\mathbb Z_3$, $\mathbb Z_3^T$), with no continuous input.\\
\emph{Gate.} Reproduce one published best fit: the weight-2 level-3 forms $Y_i(\tau)$, the mass
matrices, and the angles at the published $\tau$ (item \#12).\\
\emph{Question.} The angles at each fixed point: near-maximal $\theta_{23}$, right $\theta_{12}$
and $\theta_{13}$, which octant. \emph{Null.} The angle distribution over random $\tau$ in the
fundamental domain.\\
\emph{Result.} [pending]\quad\emph{Verdict.} [pending]

\subsection{D11: an extremum principle}\label{sec:D11}
\emph{Candidate.} Maximal CP violation:
$J=s_{12}c_{12}s_{13}c_{13}^2s_{23}c_{23}\sin\delta$ is extremal in $\theta_{23}$ at $45^\circ$ and
in $\delta$ at $\pm90^\circ$, so a maximal-$J$ principle gives the $\mu$--$\tau$-reflection pattern
without a symmetry.\\
% PENDING: theorist item #13. The proposition and its proof (first-order deviation zero;
% second-order form; other invariants at the data point).
\emph{Proposition.} [pending]\\
\emph{Question (lab).} With $\theta_{12}$ and $\theta_{13}$ at data, what deviation from $45^\circ$
does the principle allow once $\delta$ is at its data value; and is any other natural invariant
extremal at the data point? \emph{Null.} Random points.\\
\emph{Result.} [pending]\quad\emph{Verdict.} [pending]

\subsection{Summary}\label{sec:r2-summary}
% Table: candidate | near-maximal? | octant | joint prediction | (b) | (e) | lab topic | verifier check.
[pending]
